\section{Related Work}
\label{sec: related_work}

\subsection{VLM and VLA for Autonomous Driving}
\label{sec: related_work_vlm}
The extensive world knowledge and reasoning capabilities of vision-language models (VLMs) have driven their applications in autonomous driving~\cite{lmdrive, drivemm, drivegpt4, gpt-driver}. 
Early works primarily leveraged VLMs for high-level driving scene understanding and reasoning, but could not output drivable trajectories~\cite{dolphins, drivemm, drivelm, nuscenes-qa, nuinstruct, nuplanqa, lingoqa, alphadrive}. 
Subsequent methods attempted to have VLMs directly predict textual waypoints~\cite{omnidrive, hint-ad, emma, openemma}, but they struggled due to the inherent limitations of LLMs in precise numerical reasoning~\cite{math-chatgpt, mathbert}. 
This led to the development of VLA models, which integrate a planning module for end-to-end trajectory prediction~\cite{senna, opendrivevla, orion}. 
Common planning approaches include autoregressive prediction of discretized waypoints~\cite{opendrivevla, diffvla, autovla}, diffusion-based trajectory generation~\cite{drivemoe, orion, recogdrive}, and direct regression via an MLP head~\cite{simlingo}. 
However, these models suffer from sparse action supervision and often rely on auxiliary understanding and reasoning tasks to guide the learning process~\cite{opendrivevla, orion, autovla}. 
In contrast, \model employs world modeling to provide a dense signal to enhance instruction-based planning.

\subsection{World Models for Autonomous Driving}
\label{sec: related_work_world_model}
World models are typically defined as generative models that predict future states conditioned on past observations and current actions~\cite{world-models}. 
In autonomous driving, applications of world models can be categorized into three main approaches: image-based, occupancy-based, and VLA-based methods. Image-based methods leverage powerful generative architectures to synthesize high-fidelity driving videos, primarily for data generation and scene simulation~\cite{gaia-1, gaia-2, drivedreamer-1, drivedreamer-2, vista, drive-wm}. 
Occupancy-based methods model scene evolution in 3D occupancy space to enhance scene understanding~\cite{occllama, driveworld, gaussianworld} and planning~\cite{occworld, occllama, drive-occworld, preworld}, but their reliance on dense 3D labels limits scalability. 
Recently, VLA-based methods have emerged with Doe-1~\cite{doe-1} first proposing a closed-loop driving model that unifies scene understanding, prediction, and planning. 
DriveVLA-W0~\cite{drivevla-w0} integrated world modeling into a VLA framework to provide dense supervision and enhance planning. 
However, they can not perform instruction-based prediction and planning. 
Our work enables this capability, allowing the model to predict corresponding future scenes and driving trajectories conditioned on flexible language instructions.

\subsection{Unified Visual Understanding and Generation}
\label{sec: related_work_unifed}
Unified visual understanding and generation methods can be categorized into three main pipelines: quantized autoregressive (AR), external diffusion, and integrated transformers. 
Quantized AR models quantize images into discrete tokens~\cite{argen, vqgen}, enabling generation within the native autoregressive framework~\cite{janus, janus-pro, unified-io-2, tokenflow, liquid, vila-u, chameleon, emu3}. 
While this design is straightforward, its visual quality typically lags behind that of diffusion-based methods. 
The External Diffuser approach pairs a VLM with an external diffusion model~\cite{dreamllm, seed-x, emu2, metamorph}. 
The VLM provides a high-level understanding by generating a few latent tokens that condition the diffusion generator. 
However, this narrow interface between understanding and generation can restrict information flow~\cite{bagel}. 
Integrated transformer models merge autoregressive and diffusion mechanisms into a single transformer \cite{janusflow, llamafusion, transfusion_diffuse, bagel, MoT}, enabling a deep integration of powerful understanding and generation capabilities. 
In this paper, we adopt the integrated transformer to achieve instruction-based joint visual generation and action planning.
